\chapter{Implementation Details and Repository Structure}
\label{app:implementation}

\section{Repository Structure}

The project repository is organised as follows:

\begin{verbatim}
RAG Agent/
├── src/
│   ├── retriever/
│   │   ├── embeddings.py        # Embedding model wrappers (local + API)
│   │   └── retriever.py         # FAISS-based document retriever
│   ├── evaluation_agent/
│   │   ├── nli_verifier.py      # BART-MNLI NLI inference
│   │   ├── poison_detector.py   # Five-signal poison detection
│   │   ├── trust_index.py       # Trust Index computation + dampener
│   │   └── evaluation_agent.py  # Orchestration class
│   ├── generator/
│   │   └── llm_generator.py     # OpenAI-compatible LLM interface
│   └── pipeline/
│       ├── rag_pipeline.py      # End-to-end RAG pipeline
│       └── experiment_runner.py # Experiment execution + metrics
├── src/experiments/
│   └── poisoned_dataset.py      # Adversarial dataset generator
├── configs/
│   └── config.yaml              # Master configuration file
├── data/
│   ├── raw/                     # Downloaded benchmark datasets
│   ├── processed/               # Preprocessed knowledge bases
│   └── experiments/             # Serialised experiment results (JSON)
├── Master_Thesis/
│   ├── main.tex                 # LaTeX document root
│   └── chapters/                # Chapter and appendix .tex files
├── run_experiment.py            # CLI experiment runner
└── requirements.txt             # Python package dependencies
\end{verbatim}

\section{Key Dependencies}

\begin{table}[h]
\centering
\caption{Core Python dependencies and versions.}
\label{tab:app_deps}
\begin{tabular}{lll}
\hline
\textbf{Package} & \textbf{Version} & \textbf{Purpose} \\
\hline
Python         & 3.10+    & Runtime language \\
LangChain      & 0.1+     & Pipeline orchestration \\
faiss-cpu      & 1.7+     & Vector similarity search \\
transformers   & 4.35+    & BART-MNLI NLI model \\
sentence-transformers & 2.2+ & MiniLM-L6-v2 embeddings \\
torch          & 2.0+     & NLI model inference (CPU) \\
openai         & 1.0+     & FARMI API client \\
datasets       & 2.14+    & HuggingFace dataset loading \\
loguru         & 0.7+     & Structured logging \\
\hline
\end{tabular}
\end{table}

\section{Reproduction Instructions}

To reproduce the primary experiment results (Llama~3.3~70B + all-MiniLM-L6-v2, TruthfulQA, 100 samples, mixed strategy, K=3):

\begin{enumerate}
    \item Install dependencies: \texttt{pip install -r requirements.txt}
    \item Configure FARMI API credentials in \texttt{.env}: set \texttt{FARMI\_API\_KEY} and \texttt{FARMI\_API\_BASE}
    \item Run: \texttt{python run\_experiment.py}; at the interactive prompt, select Llama~3.3~70B, all-MiniLM-L6-v2, and K=3
    \item Results are saved to \texttt{data/experiments/} as a timestamped JSON file
\end{enumerate}

To reproduce the full 2×2 grid (K=3):

\begin{verbatim}
python run_experiment.py --grid --samples 50
\end{verbatim}

To reproduce the K=5 sensitivity analysis:

\begin{verbatim}
python run_experiment.py --grid --samples 50
# Select K=5 at the interactive prompt, or set TOP_K_RETRIEVAL=5 in .env
\end{verbatim}

All experiment results used in this thesis are archived in \texttt{data/experiments/} with ISO 8601 timestamps in the filename. The primary results file is:\\
\texttt{experiment\_truthfulqa\_llama3.3-70b\_all-MiniLM-L6-v2\_detection\_20260307\_112613.json}
